% =====================================================================
%  Thesis proposal, Appendix: Coupled Entry Dynamics and Embedding
%  Owns: apx:dyn
% =====================================================================

\section{Coupled entry dynamics and embedding}
\label{apx:dyn}

%--------------------------------------------------------------------------
\subsection{Coordinate rationale}\label{sec:coord-rationale}
%--------------------------------------------------------------------------

The coordinate choice is dictated by two competing requirements of the
interval-arithmetic pipeline:
\begin{enumerate}
  \item \textbf{Corridor constraints must be box constraints.} The entry
    corridor---bounded altitude, bounded speed, heating and load-factor
    limits---defines the region whose reachable subset we certify. In
    coordinates where these constraints are box faces, the interval hull is
    the natural bounding object and no wrapping is introduced by
    intersecting with the corridor. Spherical position~$(r)$ and
    wind-frame speed~$(V)$ make the primary corridor constraints
    $h_{\min}\le r - R_e \le h_{\max}$ and $V_{\min}\le V\le V_{\max}$
    into box constraints on states.

  \item \textbf{The slow manifold must be native.} The equilibrium-glide
    condition $\Lambda = 1$ is a codimension-1 surface in the
    $(r,V,\gamma)$ subspace---equivalently, once solved for the flight-path
    angle, a graph over~$(r,V)$. The profile decomposition---skeleton, bubbles,
    residual---is organized around this surface. If the slow manifold is a
    derived nonlinear function of the states (as it would be in Cartesian
    position/velocity), the separation into skeleton and bubbles inherits
    unnecessary wrapping at every step.
\end{enumerate}

The translational state is therefore carried in spherical-LVLH coordinates
$(r,\lambda,\phi,V,\gamma,\psi)$. Attitude is carried as a unit quaternion
$\vvec{q} = (q_0,q_1,q_2,q_3)^\tr$ relating the body frame~$\mathcal{B}$ and
the ECEF frame~$\mathcal{E}$ (its direction-cosine matrix
$\mat{R}_{\mathcal{E}\mathcal{B}} = \mat{R}(\vvec{q})$ is given in
\cref{eq:DCM})---this avoids all Euler-angle singularities and keeps the
quaternion kinematics in their simplest form.

The longitude equation
$\dot\lambda = V\cos\gamma\sin\psi/(r\cos\phi)$ is singular at the
geographic poles $\phi = \pm\pi/2$. This is a coordinate singularity, not a
physical one: $\lambda$ is cyclic, does not enter the corridor constraints
or the glide condition, and can be integrated separately. For any entry
scenario the ground-track extent is bounded by the energy budget, and the
coordinate system can be oriented so the trajectory stays away from the
poles.

The Savageau--Voit polynomial recasting of the thermal subsystem is dropped
entirely. Because the interval pipeline evaluates transcendentals natively,
the density-square-root embedding $y = \sqrt{\rho/\rho_0}$ and the
polynomial surrogate for wall temperature are unnecessary---the raw
physical variables are retained.


%--------------------------------------------------------------------------
\subsection{Frames}\label{sec:frames}
%--------------------------------------------------------------------------

\paragraph{Rotation convention.}
Throughout, $\mat{R}_{\mathcal{A}\mathcal{B}}$ denotes the rotation that
carries coordinates \emph{from} frame~$\mathcal{B}$ \emph{to}
frame~$\mathcal{A}$, i.e.\ $\vvec{v}^{\mathcal{A}} =
\mat{R}_{\mathcal{A}\mathcal{B}}\,\vvec{v}^{\mathcal{B}}$; its columns are the
axes of~$\mathcal{B}$ expressed in~$\mathcal{A}$. Each such matrix is
orthogonal, so $\mat{R}_{\mathcal{A}\mathcal{B}}^{-1} =
\mat{R}_{\mathcal{A}\mathcal{B}}^{\tr} = \mat{R}_{\mathcal{B}\mathcal{A}}$,
and subscripts compose in the usual way,
$\mat{R}_{\mathcal{A}\mathcal{B}}\mat{R}_{\mathcal{B}\mathcal{C}} =
\mat{R}_{\mathcal{A}\mathcal{C}}$.

\begin{itemize}
  \item $\mathcal{E}$ (\textbf{ECEF}): rotating with the constant
    planetary rate $\vvec{\omega}_E = \omega_E\,\hat{z}_{\mathcal{E}}$.

  \item $\mathcal{L}$ (\textbf{LVLH / NED}): at geocentric position
    $(r,\lambda,\phi)$, with basis vectors
    \begin{align*}
      \hat{e}_N &= -\sin\phi\cos\lambda\,\hat{x}_{\mathcal{E}}
                    -\sin\phi\sin\lambda\,\hat{y}_{\mathcal{E}}
                    +\cos\phi\,\hat{z}_{\mathcal{E}},\\
      \hat{e}_E &= -\sin\lambda\,\hat{x}_{\mathcal{E}}
                    +\cos\lambda\,\hat{y}_{\mathcal{E}},\\
      \hat{e}_D &= -\cos\phi\cos\lambda\,\hat{x}_{\mathcal{E}}
                    -\cos\phi\sin\lambda\,\hat{y}_{\mathcal{E}}
                    -\sin\phi\,\hat{z}_{\mathcal{E}}.
    \end{align*}
    The rotation $\mat{R}_{\mathcal{E}\mathcal{L}}$ from $\mathcal{L}$ to
    $\mathcal{E}$ (columns are the above vectors in~$\mathcal{E}$) is a
    function of $(\lambda,\phi)$ alone.

  \item $\mathcal{W}$ (\textbf{Wind frame}): $x$-axis along the
    planet-relative velocity, $z$-axis in the
    $\vvec{V}$--$\hat{e}_D$ plane pointing ``downward'' from the velocity,
    $y$-axis completing the right-handed system. Parameterized by
    flight-path angle~$\gamma$ (positive climbing) and
    heading~$\psi$ (clockwise from North):
    \[
      \hat{V} = \cos\gamma\cos\psi\,\hat{e}_N
              + \cos\gamma\sin\psi\,\hat{e}_E
              - \sin\gamma\,\hat{e}_D.
    \]

  \item $\mathcal{B}$ (\textbf{Body frame}): fixed at the instantaneous
    centre of mass. Related to~$\mathcal{E}$ by the attitude quaternion
    $\vvec{q}$ through the direction-cosine matrix
    $\mat{R}_{\mathcal{E}\mathcal{B}} = \mat{R}(\vvec{q})$, which carries
    body components to~$\mathcal{E}$
    ($\vvec{v}^{\mathcal{E}} = \mat{R}(\vvec{q})\,\vvec{v}^{\mathcal{B}}$):
\end{itemize}
\begin{equation}\label{eq:DCM}
  \mat{R}(\vvec{q})
  = \begin{pmatrix}
      q_0^2+q_1^2-q_2^2-q_3^2 & 2(q_1q_2-q_0q_3) & 2(q_1q_3+q_0q_2)\\
      2(q_1q_2+q_0q_3) & q_0^2-q_1^2+q_2^2-q_3^2 & 2(q_2q_3-q_0q_1)\\
      2(q_1q_3-q_0q_2) & 2(q_2q_3+q_0q_1) & q_0^2-q_1^2-q_2^2+q_3^2
    \end{pmatrix},
\end{equation}
quadratic in~$\vvec{q}$ (its columns are the body axes expressed
in~$\mathcal{E}$).


%--------------------------------------------------------------------------
\subsection{State vector}\label{sec:state-vector}
%--------------------------------------------------------------------------

The full coupled state is
\begin{equation}\label{eq:state-full}
  x = \bigl(\underbrace{r,\lambda,\phi,V,\gamma,\psi}_{6},\;
       \underbrace{q_0,q_1,q_2,q_3,p,q_b,r_b}_{7},\;
       \underbrace{Q_{\mathrm{tot}},T_w,s,m}_{4},\;
       \underbrace{\eta_1,\dots,\eta_{N_m},
                   \dot\eta_1,\dots,\dot\eta_{N_m}}_{2N_m}\bigr),
\end{equation}
with the unit-quaternion constraint $\norm{\vvec{q}}=1$ making the
translational-rotational block effectively 12-dimensional. Total dimension:
$17 + 2N_m$.


%--------------------------------------------------------------------------
\subsection{Translational kinematics}\label{sec:trans-kin}
%--------------------------------------------------------------------------

The planet-relative velocity in the $\mathcal{L}$ (NED) frame is
$\vvec{V}_{\mathcal{L}} =
(V\cos\gamma\cos\psi,\;V\cos\gamma\sin\psi,\;-V\sin\gamma)^\tr$. The
kinematic equations are:
\begin{align}
  \dot r      &= V\sin\gamma, \label{eq:rdot}\\
  \dot\lambda &= \frac{V\cos\gamma\sin\psi}{r\cos\phi},
                                                         \label{eq:lamdot}\\
  \dot\phi    &= \frac{V\cos\gamma\cos\psi}{r}.         \label{eq:phidot}
\end{align}


%--------------------------------------------------------------------------
\subsection{Gravitational acceleration}\label{sec:gravity}
%--------------------------------------------------------------------------

The $J_2$ gravitational potential in geocentric spherical coordinates is
\[
  \Phi_g(r,\phi)
  = -\frac{\mu}{r}\biggl[
      1 - \half J_2\Bigl(\frac{R_e}{r}\Bigr)^{\!2}
          \bigl(3\sin^2\!\phi - 1\bigr)
    \biggr].
\]
Its gradient, resolved in $\mathcal{L}$ (NED), has only radial and
northward components:
\begin{align}
  g_D(r,\phi) &= \frac{\mu}{r^2}\biggl[
    1 + \tfrac{3}{2}\,J_2\Bigl(\frac{R_e}{r}\Bigr)^{\!2}
        \bigl(1-5\sin^2\!\phi\bigr)
  \biggr], \label{eq:gD}\\[4pt]
  g_N(r,\phi) &= -\frac{3\mu}{r^2}\,J_2
    \Bigl(\frac{R_e}{r}\Bigr)^{\!2}
    \sin\phi\cos\phi, \label{eq:gN}\\
  g_E &= 0\quad\text{(by axial symmetry)}.\nonumber
\end{align}

\begin{remark}\label{rem:gravity-spherical}
This reproduces the Cartesian form under the
identifications $r_3/\norm{\vvec{r}} = \sin\phi$ and
$\norm{\vvec{r}} = r$. The spherical form makes the interval evaluation
tighter: $g_D$ and~$g_N$ are explicit functions of $(r,\phi)$, with no
square root of a sum of squares.
\end{remark}


%--------------------------------------------------------------------------
\subsection{Translational dynamics}\label{sec:trans-dyn}
%--------------------------------------------------------------------------

Newton's law in the rotating ECEF frame, resolved along the wind-frame
axes (tangential~$\hat{V}$, normal-in-vertical-plane, lateral), gives
the force equations. Define the aerodynamic force components:
\begin{align*}
  F_V     &= \text{force along } \hat{V},\\
  F_\gamma &= \text{force $\perp\hat{V}$, in the vertical plane},\\
  F_\psi  &= \text{force $\perp$ the vertical plane (lateral)}.
\end{align*}
In the 3-DOF model $F_V=-D$, $F_\gamma=L\cos\sigma$,
$F_\psi=L\sin\sigma$. In the full 6-DOF formulation these are
\emph{derived from the attitude quaternion}---see
\cref{sec:aero-force}.

The translational dynamical equations are:
\begin{align}
  \dot V &= \frac{F_V}{m}
    - g_D\sin\gamma + g_N\cos\gamma\cos\psi
    + \omega_E^2 r\cos\phi
      \bigl[\sin\gamma\cos\phi - \cos\gamma\cos\psi\sin\phi\bigr],
    \label{eq:Vdot}
  \\[6pt]
  \dot\gamma &= \frac{1}{V}\biggl\{
    \frac{F_\gamma}{m}
    - \Bigl(g_D - \frac{V^2}{r}\Bigr)\cos\gamma
    - g_N\sin\gamma\cos\psi
    + 2\omega_E V\cos\phi\sin\psi \nonumber\\
    &\qquad{}+ \omega_E^2 r\cos\phi
      \bigl[\cos\gamma\cos\phi + \sin\gamma\cos\psi\sin\phi\bigr]
    \biggr\}, \label{eq:gamdot}
  \\[6pt]
  \dot\psi &= \frac{1}{V\cos\gamma}\biggl\{
    \frac{F_\psi}{m}
    - \frac{V^2}{r}\cos\gamma\sin\psi\tan\phi
    + g_N\sin\psi \nonumber\\
    &\qquad{}+ 2\omega_E V
      \bigl[\tan\gamma\cos\phi\cos\psi - \sin\phi\bigr]
    + \omega_E^2 r\cos\phi\sin\phi
      \frac{\sin\psi}{\cos\gamma}
    \biggr\}. \label{eq:psidot}
\end{align}

\begin{remark}\label{rem:vinh-reduction}
With $g_N = 0$ (spherical gravity) and
$\omega_E = 0$ (non-rotating planet), these reduce to the classical Vinh
equations for unpowered entry~\cite{vinh1980hypersonic}.
\end{remark}

\begin{remark}\label{rem:psi-singularity}
The singularity $1/\cos\gamma$ in~$\dot\psi$ occurs
at vertical flight ($\gamma = \pm\pi/2$). This is not encountered during
the equilibrium-glide phase and occurs at most instantaneously during a
skip peak. For the interval pipeline, the singular factor is enclosed by
standard interval evaluation of $1/\cos[\gamma]$ provided the interval
$[\gamma]$ does not contain $\pm\pi/2$.
\end{remark}


%--------------------------------------------------------------------------
\subsection{Rotational dynamics}\label{sec:rot-dyn}
%--------------------------------------------------------------------------

The quaternion kinematics, driven by the body rate relative
to~$\mathcal{E}$:
\begin{equation}\label{eq:qdot}
  \dot{\vvec{q}} = \half\,\mat{\Omega}(\vvec{\omega})\,\vvec{q},
\end{equation}
with $\vvec{\omega} = (p,q_b,r_b)^\tr$ the inertial angular velocity in
body-frame components and
\begin{equation}\label{eq:Omega}
  \mat{\Omega}(\vvec{a})
  = \begin{pmatrix}
      0    & -a_1 & -a_2 & -a_3\\
      a_1  &  0   &  a_3 & -a_2\\
      a_2  & -a_3 &  0   &  a_1\\
      a_3  &  a_2 & -a_1 &  0
    \end{pmatrix}.
\end{equation}

Euler's equations, with inertia tensor
$\mat{I} = \mat{I}(s,\eta)$ depending on recession depth and modal
amplitudes:
\begin{equation}\label{eq:euler}
  \mat{I}\,\dot{\vvec{\omega}}
  = \vvec{M} - \vvec{\omega}\times(\mat{I}\,\vvec{\omega})
    - \dot{\mat{I}}\,\vvec{\omega},
\end{equation}
where $\vvec{M} = \vvec{M}_a + \vvec{M}_{\mathrm{rcs}}$ is the total
external moment in~$\mathcal{B}$.


%--------------------------------------------------------------------------
\subsection{Aerodynamic force: from quaternion to
  \texorpdfstring{$(F_V,F_\gamma,F_\psi)$}
  {(FV, Fgamma, Fpsi)}}\label{sec:aero-force}
%--------------------------------------------------------------------------

This section is the bridge between the rotational and translational
dynamics.

\paragraph{Step 1: Body-frame aerodynamic force.}
Given angle of attack~$\alpha$, sideslip~$\beta$, dynamic pressure
$q_{\mathrm{dyn}} = \half\rho V^2$, and reference area~$S$:
\begin{equation}\label{eq:Fa-body}
  \vvec{F}_a^{\mathcal{B}}
  = q_{\mathrm{dyn}}\,S\,
    \bigl(-C_A(\alpha,\beta,M),\;
           C_Y(\alpha,\beta,M),\;
          -C_N(\alpha,\beta,M)\bigr)^\tr.
\end{equation}

\paragraph{Step 2: Transform to LVLH.}
Define the body-to-LVLH rotation
\begin{equation}\label{eq:RLB}
  \mat{R}_{\mathcal{L}\mathcal{B}}({\vvec{q}},\lambda,\phi)
  \;:=\;
  \mat{R}_{\mathcal{E}\mathcal{L}}^\tr(\lambda,\phi)\;
  \mat{R}(\vvec{q})
  \;=\;
  \mat{R}_{\mathcal{L}\mathcal{E}}\,\mat{R}_{\mathcal{E}\mathcal{B}},
\end{equation}
so that $\vvec{v}^{\mathcal{L}} =
\mat{R}_{\mathcal{L}\mathcal{B}}\,\vvec{v}^{\mathcal{B}}$. Then the
aerodynamic force in~$\mathcal{L}$ is
\begin{equation}\label{eq:Fa-LVLH}
  \vvec{F}_a^{\mathcal{L}}
  = \mat{R}_{\mathcal{L}\mathcal{B}}\,
    \vvec{F}_a^{\mathcal{B}}.
\end{equation}

\paragraph{Step 3: Project onto wind axes.}
The wind-axis triad, expressed in~$\mathcal{L}$, is
\begin{align}
  \hat{V}_{\mathcal{L}} &= (\cos\gamma\cos\psi,\;
                      \cos\gamma\sin\psi,\;-\sin\gamma)^\tr,
                                                    \label{eq:Vhat}\\
  \hat{n}_\gamma &= (-\sin\gamma\cos\psi,\;
                      -\sin\gamma\sin\psi,\;
                      -\cos\gamma)^\tr, \label{eq:ngam}\\
  \hat{n}_\psi   &= \hat{V}_{\mathcal{L}} \times \hat{n}_\gamma.
                                                    \label{eq:npsi}
\end{align}
Then:
\begin{equation}\label{eq:F-projections}
  F_V = \vvec{F}_a^{\mathcal{L}}\cdot\hat{V}_{\mathcal{L}},\qquad
  F_\gamma = \vvec{F}_a^{\mathcal{L}}\cdot\hat{n}_\gamma,\qquad
  F_\psi = \vvec{F}_a^{\mathcal{L}}\cdot\hat{n}_\psi.
\end{equation}

\paragraph{Step 4: Identification with lift, drag, bank angle.}
In the point-mass limit where the attitude is in trim,
$F_V = -D$, $F_\gamma = L\cos\sigma$, $F_\psi = L\sin\sigma$.
In the full 6-DOF formulation $D$, $L$, $\sigma$ are \emph{not independent
inputs}: they are determined by the attitude quaternion~$\vvec{q}$, the
aerodynamic coefficients, and the velocity direction $(\gamma,\psi)$
through Eqs.~(\ref{eq:Fa-body})--(\ref{eq:F-projections}).

\begin{remark}\label{rem:bank-angle}
The effective bank angle can be extracted as
$\sigma = \operatorname{atan2}(F_\psi, F_\gamma)$,
an algebraic function of
$({\vvec{q}},\lambda,\phi,\gamma,\psi,\alpha,\beta)$. Under interval
arithmetic, $\sigma$ is evaluated as an interval $\operatorname{atan2}$, a
standard operation in validated numerics~\cite{moore2009introduction}.
\end{remark}

\begin{remark}\label{rem:aoa-sideslip}
The angle of attack~$\alpha$ and sideslip~$\beta$ are
determined by the quaternion and velocity direction. Let
$\vvec{V}^{\mathcal{B}} = \mat{R}_{\mathcal{L}\mathcal{B}}^{\tr}\,
\vvec{V}_{\mathcal{L}}$ be the velocity in body-frame components. Then
$\alpha = \operatorname{atan2}(-V_z^{\mathcal{B}}, -V_x^{\mathcal{B}})$
and $\beta = \arcsin(V_y^{\mathcal{B}}/V)$.
\end{remark}


%--------------------------------------------------------------------------
\subsection{Equilibrium-glide slow manifold}\label{sec:glide-manifold}
%--------------------------------------------------------------------------

The equilibrium-glide condition is the slow manifold of the translational
dynamics under the thin-atmosphere singular perturbation
$\eps_{\mathrm{atm}} = H_s/R_e \to 0$.

Setting $\dot\gamma = 0$ and neglecting Earth rotation ($\omega_E = 0$) in
Eq.~(\ref{eq:gamdot}):
\[
  \frac{F_\gamma}{m}
  = \Bigl(g_D - \frac{V^2}{r}\Bigr)\cos\gamma
    + g_N\sin\gamma\cos\psi.
\]
For small flight-path angle ($\gamma\ll 1$) and spherical gravity
($g_N = 0$, $g_D = \mu/r^2$):
\begin{equation}\label{eq:glide-3dof}
  \frac{F_\gamma}{m} = g - \frac{V^2}{r},
  \qquad\text{i.e.}\quad
  \frac{L\cos\sigma}{m(g - V^2/r)} = 1
  \quad(\Lambda = 1).
\end{equation}

In the full 6-DOF formulation, $F_\gamma$ is determined by the quaternion
via Eqs.~(\ref{eq:Fa-body})--(\ref{eq:F-projections}). The slow-manifold
condition therefore constrains the quaternion directly:
\begin{multline}\label{eq:glide-full}
  \bigl[\mat{R}_{\mathcal{L}\mathcal{B}}\,
         \vvec{F}_a^{\mathcal{B}}\bigr]
  \cdot\hat{n}_\gamma \\
  = m\Bigl(g_D(r,\phi) - \frac{V^2}{r}\Bigr)\cos\gamma
    + m\,g_N(r,\phi)\sin\gamma\cos\psi,
\end{multline}
a single scalar constraint on the coupled translational and rotational state
(dimension~$13$); it defines the codimension-1 slow manifold~$\Lambda$.

\begin{remark}\label{rem:manifold-attitude}
The slow manifold is \emph{not} a constraint on
$\sigma$ alone---it constrains a specific component of
$\mat{R}_{\mathcal{L}\mathcal{B}} \vvec{F}_a^{\mathcal{B}}$, which
depends on the full quaternion and on the aerodynamic angles
$(\alpha,\beta)$ implicitly. The ``bank angle'' that appears in the
point-mass glide condition is a derived quantity, and the slow manifold
constrains the attitude directly.
\end{remark}

\begin{remark}\label{rem:j2-glide}
Equation~(\ref{eq:glide-full}) retains the $J_2$
contributions through $g_D(r,\phi)$ and $g_N(r,\phi)$. The simplified
form $\Lambda = L\cos\sigma/[m(g-V^2/r)] = 1$ is the leading-order
reduction under spherical gravity ($J_2 = 0$) and small flight-path angle.
The $J_2$ correction to the glide balance is $O(J_2)$, ranging from
${\approx}\,0.16\%$ of the leading gravitational term near the equator to
${\approx}\,0.65\%$ near the poles (through the $(1-5\sin^2\!\phi)$ factor);
for the interval pipeline it is enclosed exactly.
\end{remark}

\begin{remark}\label{rem:j2-skip}
The $J_2$ contribution is especially significant
during skip maneuvers. When the vehicle leaves the slow manifold (the
defining condition of a boundary-layer bubble), three $J_2$ effects
compound:
\begin{enumerate}
  \item The northward gravitational component $g_N(r,\phi)$ enters the
    heading rate $\dot\psi$ (Eq.~\ref{eq:psidot}), producing a
    latitude-dependent heading drift during each skip that is absent under
    spherical gravity. Northward and southward skips are not symmetric.
  \item The radial component $g_D(r,\phi)$ varies with latitude by up to
    $\approx 0.65\%$ of the leading term through the $(1-5\sin^2\!\phi)$
    factor, modifying the
    effective centrifugal balance $g_D - V^2/r$ that governs the
    re-entry angle after each skip peak.
  \item Over $K$ skips the ground track may span a significant latitude
    range, so the $J_2$ corrections at successive re-entries are evaluated
    at different~$\phi$, and the effects are cumulative.
\end{enumerate}
For the interval pipeline, these effects are enclosed automatically: the
interval evaluations of $g_D([r],[\phi])$ and $g_N([r],[\phi])$ widen with
the latitude interval, which is itself widest during a skip (where the
ground track sweeps the most latitude). Retaining $J_2$ in the dynamics
is therefore not a refinement but a necessity for tight certified bounds on
skip trajectories.
\end{remark}


%--------------------------------------------------------------------------
\subsection{Thermal subsystem (native form)}\label{sec:thermal}
%--------------------------------------------------------------------------

The atmosphere density profile is $\rho(h)$ from the NRLMSIS~2.0
model~\cite{emmert2021nrlmsis}, with $h = r - R_e$. Under interval arithmetic,
$\rho(h)$ is enclosed by evaluating the model's analytic approximation (or
a validated piecewise-exponential envelope) with interval
argument~$[h]$.

Stagnation-point heating (Sutton--Graves~\cite{sutton1971general}, within its
stated validity envelope: axisymmetric blunt bodies, chemical equilibrium,
enthalpy 2.3--116.2\,MJ/kg, pressure 0.001--100\,atm):
\begin{equation}\label{eq:Qdot}
  \dot Q_{\mathrm{tot}} = k_Q\sqrt{\rho/R_n}\,V^3.
\end{equation}
Wall temperature:
\begin{equation}\label{eq:Twdot}
  \dot T_w
  = \frac{\dot Q_{\mathrm{tot}} - (T_w - T_\infty)/\tau_{\mathrm{cond}}}
         {C_w}.
\end{equation}
No embedding, no recasting. The transcendentals $\sqrt{\rho}$ and $V^3$
are evaluated directly as interval operations.

\begin{remark}\label{rem:thermal-embedding}
The density-square-root embedding
$y = \sqrt{\rho/\rho_0}$ of the original formulation was introduced to
render the heating correlation polynomial in the states (the
Savageau--Voit recasting strategy~\cite{savageau1987recasting}). Since the interval
pipeline evaluates $\sqrt{\cdot}$ natively, this auxiliary state is
eliminated, and with it the associated wrapping from carrying an extra
dimension.
\end{remark}


%--------------------------------------------------------------------------
\subsection{Ablation}\label{sec:ablation}
%--------------------------------------------------------------------------

Recession rate coupled to heat flux; mass loss from char-layer blowoff:
\begin{equation}\label{eq:ablation}
  \dot s = \alpha_s\,\dot Q_{\mathrm{tot}},\qquad
  \dot m = -\dot m_{\mathrm{abl}}.
\end{equation}


%--------------------------------------------------------------------------
\subsection{Modal / structural elasticity}\label{sec:modal}
%--------------------------------------------------------------------------

Per mode~$i$, with piston-theory forcing and temperature-dependent natural
frequency:
\begin{equation}\label{eq:modal}
  \ddot\eta_i + 2\zeta_i\,\omega_i(T_w)\,\dot\eta_i
  + \omega_i(T_w)^2\,\eta_i
  = Q_i(x,\eta,\dot\eta).
\end{equation}
Piston theory is a double expansion in thickness ratio and inverse Mach
number valid when the contributions of order $\delta/M^3$ and
$\delta^2/M^2$ are negligible~\cite{ashley1956piston,lighthill1953oscillating}; its
applicability conditions must be checked at each atmospheric pass. At
skip peaks with high angle of attack and strong shock interaction, piston
theory may fail, and the modal forcing should be flagged as uncertain
(widened interval) in those regimes.


%--------------------------------------------------------------------------
\subsection{Density model and interval envelope}\label{sec:density}
%--------------------------------------------------------------------------

\emph{To be completed:} functional form of the atmospheric density profile
suitable for interval evaluation. The NRLMSIS~2.0 model provides the
reference; the interval pipeline requires either a closed-form analytic
approximation with certified error bounds, or a piecewise-exponential
envelope
\[
  \rho_{\mathrm{lower}}(h) \le \rho(h) \le \rho_{\mathrm{upper}}(h)
\]
valid across the altitude range of the corridor. The width
$\rho_{\mathrm{upper}} - \rho_{\mathrm{lower}}$ at each altitude is the
primary driver of the aerodynamic interval widths.


%--------------------------------------------------------------------------
\subsection{Summary of coordinate choices}\label{sec:summary-table}
%--------------------------------------------------------------------------

\begin{center}
\begin{tabular}{@{}lll@{}}
\toprule
\textbf{Feature} &
\textbf{Spherical-LVLH + Quat.} &
\textbf{Cartesian ECEF + Quat.} \\
\midrule
Corridor constraints &
  Box in $(r,V)$ &
  Nonlinear: $h = \norm{\vvec{r}} - R_e$ \\
Glide condition &
  Native: Eq.~(\ref{eq:glide-3dof})/(\ref{eq:glide-full}) &
  Derived from $(u,v,w)$, $\norm{\vvec{r}}$ \\
Position singularity &
  Pole in $\dot\lambda$ &
  None \\
Dynamics &
  Transcendental ($\sin$, $\cos$) &
  Polynomial (quadratic in $\vvec{q}$) \\
Wrapping source &
  Trig dependency &
  Corridor intersection \\
Density evaluation &
  $\rho(r-R_e)$: direct &
  $\rho(\norm{\vvec{r}}-R_e)$: extra $\sqrt{\cdot}$ \\
\bottomrule
\end{tabular}
\end{center}

\smallskip\noindent
The spherical-LVLH choice concentrates wrapping in the dynamics
(trigonometric evaluations), where it is distributed across every time step
and can be controlled by step-size refinement. The Cartesian choice
concentrates wrapping in the corridor intersection, where it is applied
once per intersection and inflates the entire reachable-set bound. For the
profile decomposition---which is organized around the corridor and the
glide manifold---the spherical choice produces tighter certified bounds.
